\subsubsection{Spectral Theory}

Spectral theory studies an operator through algebraic expressions in the operator and through bases that reveal invariant flags. Polynomial expressions produce eigenvalues over $\mathbb{C}$, and eigenvalues make upper-triangular representations possible.

\begin{definition} \label{def:operator-powers} \hlt{Powers of an Operator}\\
Let $T\in\LL(\mathcal{V})$. For a positive integer $m$, $T^m\defeq \underbrace{T\cdots T}_{m\text{ times}}$ and $T^0\defeq I_{\mathcal{V}}$.\\
If $T$ is invertible, then $T^{-m}\defeq(T^{-1})^m$ for positive integers $m$.
\end{definition}

\begin{definition} \label{def:polynomial-applied-to-operator} \hlt{Polynomial Applied to an Operator}\\
Let $T\in\LL(\mathcal{V})$ and let $p\in\mathcal{P}(\mathbb{F})$ be $p(z)=a_0+a_1z+\cdots+a_mz^m$. Then $p(T)\defeq a_0I_{\mathcal{V}}+a_1T+\cdots+a_mT^m$.
\end{definition}

\begin{definition} \label{def:product-of-polynomials} \hlt{Product of Polynomials}\\
If $p,q\in\mathcal{P}(\mathbb{F})$, their product $pq\in\mathcal{P}(\mathbb{F})$ is defined by $(pq)(z)=p(z)q(z)$ for $z\in\mathbb{F}$.
\end{definition}

\begin{proposition} \label{prop:operator-polynomial-products} \hlt{Multiplicative Properties of Operator Polynomials}\\
Let $p,q\in\mathcal{P}(\mathbb{F})$ and $T\in\LL(\mathcal{V})$. Then
\begin{enumerate}[label=(\roman*)]
\setlength{\itemsep}{0pt}
\item $(pq)(T)=p(T)q(T)$;
\item $p(T)q(T)=q(T)p(T)$.
\end{enumerate}
\begin{proof}
\begin{enumerate}[label=(\roman*)]
\setlength{\itemsep}{0pt}
\item Write $p(z)=\sum_{j=0}^{m}a_jz^j$ and $q(z)=\sum_{k=0}^{n}b_kz^k$. Then $(pq)(z)=\sum\limits_{j=0}^{m}\sum\limits_{k=0}^{n}a_jb_kz^{j+k}$.\\
Substituting $T$ gives $(pq)(T)=\sum\limits_{j=0}^{m}\sum\limits_{k=0}^{n}a_jb_kT^{j+k} =\left(\sum\limits_{j=0}^{m}a_jT^j\right)\left(\sum\limits_{k=0}^{n}b_kT^k\right) =p(T)q(T)$.
\item Since polynomial multiplication is commutative, $pq=qp$.\\
By part (i), $p(T)q(T)=(pq)(T)=(qp)(T)=q(T)p(T)$.
\end{enumerate}
\hlt{Structural Commentary}: Polynomial evaluation preserves products.\\
\hlt{Key Insight}: Expanding a polynomial product works the same with $T$ as with $z$.
\end{proof}
\end{proposition}

\begin{proposition} \label{prop:kernel-of-polynomial-in-operator-is-invariant} \hlt{Kernel of a Polynomial in an Operator Is Invariant}\\
Let $T\in\LL(\mathcal{V})$ and $p\in\mathcal{P}(\mathbb{F})$. Then $\ker(p(T))$ is invariant under $T$.
\begin{proof}
If $\mathbf{v}\in\ker(p(T))$, then $p(T)\mathbf{v}=\mathbf{0}$. Since $p(T)$ commutes with $T$ by Proposition \ref{prop:operator-polynomial-products}, $p(T)(T\mathbf{v})=T(p(T)\mathbf{v})=T\mathbf{0}=\mathbf{0}$.\\
\hlt{Structural Commentary}: Polynomial kernels are $T$-invariant.\\
\hlt{Key Insight}: Applying $T$ after a polynomial annihilation preserves the annihilation.
\end{proof}
\end{proposition}

\begin{theorem} \label{thm:complex-operator-has-eigenvalue} \hlt{Complex Operators Have Eigenvalues}\\
Every operator on a finite-dimensional, nonzero complex vector space has an eigenvalue.
\begin{proof}
Let $\mathcal{V}$ be a complex vector space with $\dim(\mathcal{V})=n>0$, and let $T\in\LL(\mathcal{V})$.\\
Choose $\mathbf{v}\neq\mathbf{0}$. The list $(\mathbf{v},T\mathbf{v},T^2\mathbf{v},\ldots,T^n\mathbf{v})$ has $n+1$ vectors in an $n$-dimensional space, so it is linearly dependent. Hence for some scalars $a_0,\ldots,a_n\in\mathbb{C}$, not all zero, $a_0\mathbf{v}+a_1T\mathbf{v}+\cdots+a_nT^n\mathbf{v}=\mathbf{0}$.\\
Let $m$ be the largest index with $a_m\neq0$. Since $\mathbf{v}\neq\mathbf{0}$, not all of $a_1,\ldots,a_n$ can be zero, so $m\geq1$. Set $p(z)=a_0+a_1z+\cdots+a_mz^m$. Then $p(T)\mathbf{v}=\mathbf{0}$. By the Fundamental Theorem of Algebra, $p(z)=c(z-\lambda_1)\cdots(z-\lambda_m)$ for some $c\neq0$ and $\lambda_1,\ldots,\lambda_m\in\mathbb{C}$. Thus $c(T-\lambda_1I_{\mathcal{V}})\cdots(T-\lambda_mI_{\mathcal{V}})\mathbf{v}=\mathbf{0}$.\\
Because $\mathbf{v}\neq\mathbf{0}$, at least one factor $T-\lambda_jI_{\mathcal{V}}$ is not injective; otherwise their product would be injective. Hence $\ker(T-\lambda_jI_{\mathcal{V}})$ contains a nonzero vector, so $\lambda_j$ is an eigenvalue by Proposition \ref{prop:eigenvectors-as-kernel-vectors}.\\
\hlt{Structural Commentary}: A vector orbit gives a splitting annihilating polynomial.\\
\hlt{Key Insight}: A product with nontrivial kernel has a nontrivial factor.
\end{proof}
\end{theorem}

\begin{definition} \label{def:diagonal-and-upper-triangular-matrices} \hlt{Diagonal and Upper-Triangular Matrices}\\
The \hlt{diagonal} of a square matrix is its upper-left to lower-right entries. A square matrix is \hlt{diagonal} if all off-diagonal entries are zero, and \hlt{upper triangular} if all entries below the diagonal are zero.
\end{definition}

\begin{definition} \label{def:diagonalizable-operator} \hlt{Diagonalizable Operator}\\
An operator $T\in\LL(\mathcal{V})$ is \hlt{diagonalizable} if $T$ has a diagonal matrix with respect to some basis of $\mathcal{V}$.
\end{definition}

\begin{proposition} \label{prop:eigenspaces-for-distinct-eigenvalues-form-direct-sum} \hlt{Eigenspaces for Distinct Eigenvalues Form a Direct Sum}\\
Let $\mathcal{V}$ be finite-dimensional, let $T\in\LL(\mathcal{V})$, and suppose $\lambda_1,\ldots,\lambda_m$ are distinct eigenvalues of $T$.\\
Then $E_{\lambda_1}(T)+\cdots+E_{\lambda_m}(T)$ is a direct sum. Consequently, $\dim E_{\lambda_1}(T)+\cdots+\dim E_{\lambda_m}(T)\leq\dim(\mathcal{V})$.
\begin{proof}
Suppose $\mathbf{u}_1+\cdots+\mathbf{u}_m=\mathbf{0}$ with $\mathbf{u}_j\in E_{\lambda_j}(T)$. If some $\mathbf{u}_j$ were nonzero, then the nonzero vectors among $\mathbf{u}_1,\ldots,\mathbf{u}_m$ would be eigenvectors corresponding to distinct eigenvalues, and Theorem \ref{thm:distinct-eigenvalues-independent-eigenvectors} would make them linearly independent. This contradicts their sum being $\mathbf{0}$. Hence every $\mathbf{u}_j=\mathbf{0}$, so the sum is direct.\\
By Theorem \ref{thm:direct-sum-iff-dimensions-add}, $\dim(E_{\lambda_1}(T)+\cdots+E_{\lambda_m}(T))=\dim E_{\lambda_1}(T)+\cdots+\dim E_{\lambda_m}(T)$.\\
The left side is the dimension of a subspace of $\mathcal{V}$, so Proposition \ref{prop:dimension-bounds-and-basis-criteria} gives the stated inequality.\\
\hlt{Structural Commentary}: Distinct eigenspaces form independent pieces.\\
\hlt{Key Insight}: A zero sum across them would contradict eigenvector independence.
\end{proof}
\end{proposition}

\begin{theorem} \label{thm:conditions-equivalent-to-diagonalizability} \hlt{Conditions Equivalent to Diagonalizability}\\
Let $\mathcal{V}$ be finite-dimensional, let $T\in\LL(\mathcal{V})$, and let $\lambda_1,\ldots,\lambda_m$ denote the distinct eigenvalues of $T$. Set $n=\dim(\mathcal{V})$. The following are equivalent:
\begin{enumerate}[label=(\roman*)]
\setlength{\itemsep}{0pt}
\item $T$ is diagonalizable;
\item $\mathcal{V}$ has a basis consisting of eigenvectors of $T$;
\item $\exists$ one-dimensional subspaces $\mathcal{U}_1,\ldots,\mathcal{U}_n$ of $\mathcal{V}$, each invariant under $T$, such that $\mathcal{V}=\mathcal{U}_1\oplus\cdots\oplus\mathcal{U}_n$;
\item $\mathcal{V}=E_{\lambda_1}(T)\oplus\cdots\oplus E_{\lambda_m}(T)$;
\item $\dim(\mathcal{V})=\dim E_{\lambda_1}(T)+\cdots+\dim E_{\lambda_m}(T)$.
\end{enumerate}
\begin{proof}
\begin{enumerate}[label=(\roman*).]
\setlength{\itemsep}{0pt}
\item[$(i)$]$\Rightarrow (ii)$ If $T$ is diagonalizable, then $T$ has a diagonal matrix with respect to some basis $(\mathbf{v}_1,\ldots,\mathbf{v}_n)$. Hence $T\mathbf{v}_j=a_j\mathbf{v}_j$ for each $j$, so this basis consists of eigenvectors of $T$.
\item[$(ii)$]$\Rightarrow (i)$ If $(\mathbf{v}_1,\ldots,\mathbf{v}_n)$ is a basis of eigenvectors of $T$, then $T\mathbf{v}_j=a_j\mathbf{v}_j$ for each $j$. Thus the matrix of $T$ with respect to this basis is diagonal.
\item[$(ii)$]$\Rightarrow (iii)$ Let $(\mathbf{v}_1,\ldots,\mathbf{v}_n)$ be an eigenbasis, and put $\mathcal{U}_j=\operatorname{span}(\mathbf{v}_j)$. Each $\mathcal{U}_j$ is one-dimensional and invariant under $T$ by Proposition \ref{prop:one-dimensional-invariant-subspaces-and-eigenvectors}; the basis decomposition gives $\mathcal{V}=\mathcal{U}_1\oplus\cdots\oplus\mathcal{U}_n$.
\item[$(iii)$]$\Rightarrow (ii)$ Choose nonzero $\mathbf{v}_j\in\mathcal{U}_j$. Proposition \ref{prop:one-dimensional-invariant-subspaces-and-eigenvectors} makes each $\mathbf{v}_j$ an eigenvector of $T$, and the direct-sum decomposition makes $(\mathbf{v}_1,\ldots,\mathbf{v}_n)$ a basis of $\mathcal{V}$.
\item[$(ii)$]$\Rightarrow (iv)$ Every vector in an eigenbasis lies in one of $E_{\lambda_1}(T),\ldots,E_{\lambda_m}(T)$, so these eigenspaces span $\mathcal{V}$. Proposition \ref{prop:eigenspaces-for-distinct-eigenvalues-form-direct-sum} shows their sum is direct.
\item[$(iv)$]$\Rightarrow (v)$ If $\mathcal{V}=E_{\lambda_1}(T)\oplus\cdots\oplus E_{\lambda_m}(T)$, then dimensions add by Theorem \ref{thm:direct-sum-iff-dimensions-add}.
\item[$(v)$]$\Rightarrow (ii)$ For each $j$, choose a basis $\beta_j$ of $E_{\lambda_j}(T)$, and concatenate these bases into a list $\beta$. The length of $\beta$ is $\dim(\mathcal{V})$. If a linear combination of vectors in $\beta$ is zero, grouping terms by eigenspace gives $\mathbf{u}_1+\cdots+\mathbf{u}_m=\mathbf{0}$ with $\mathbf{u}_j\in E_{\lambda_j}(T)$. Proposition \ref{prop:eigenspaces-for-distinct-eigenvalues-form-direct-sum} gives $\mathbf{u}_j=\mathbf{0}$ for every $j$; since each $\beta_j$ is independent, all coefficients are zero. Thus $\beta$ is independent, Proposition \ref{prop:dimension-bounds-and-basis-criteria} makes $\beta$ a basis of $\mathcal{V}$ consisting of eigenvectors of $T$.
\end{enumerate}
\hlt{Structural Commentary}: Diagonalizability means the eigenspaces exhaust $\mathcal{V}$. Directness is automatic for distinct eigenspaces, so the real test is spanning, or equivalently the dimension count.\\
\hlt{Key Insight}: A diagonal basis is exactly an eigenbasis.
\end{proof}
\end{theorem}

\begin{perspective} \label{per:diagonalizability-viewpoint} \hlt{Diagonalization as Exhaustion by Eigenspaces}\\
Distinct eigenspaces are already separated by Proposition \ref{prop:eigenspaces-for-distinct-eigenvalues-form-direct-sum}. Thus diagonalizability is the stronger statement that these separated pieces account for every direction in $\mathcal{V}$.
\end{perspective}

\begin{corollary} \label{cor:enough-eigenvalues-implies-diagonalizable} \hlt{Enough Eigenvalues Imply Diagonalizability}\\
If $\mathcal{V}$ is finite-dimensional and $T\in\LL(\mathcal{V})$ has $\dim(\mathcal{V})$ distinct eigenvalues, then $T$ is diagonalizable.
\begin{proof}
Let $n=\dim(\mathcal{V})$, and let $\lambda_1,\ldots,\lambda_n$ be distinct eigenvalues of $T$. For each $j$, choose an eigenvector $\mathbf{v}_j$ corresponding to $\lambda_j$. By Theorem \ref{thm:distinct-eigenvalues-independent-eigenvectors}, $(\mathbf{v}_1,\ldots,\mathbf{v}_n)$ is linearly independent. Since this list has length $n$, Proposition \ref{prop:dimension-bounds-and-basis-criteria} makes it a basis of $\mathcal{V}$. Thus $\mathcal{V}$ has an eigenbasis, so $T$ is diagonalizable by Theorem \ref{thm:conditions-equivalent-to-diagonalizability}.\\
\hlt{Structural Commentary}: One eigenvector from each distinct eigenvalue already gives a full-size independent list.\\
\hlt{Key Insight}: Once the independent eigenvectors reach $\dim(\mathcal{V})$, they must form a basis.
\end{proof}
\end{corollary}

\begin{theorem} \label{thm:upper-triangular-criteria} \hlt{Criteria for an Upper-Triangular Operator Matrix}\\
Let $T\in\LL(\mathcal{V})$ and let $\beta=(\mathbf{v}_1,\ldots,\mathbf{v}_n)$ be a basis of $\mathcal{V}$. The following are equivalent:
\begin{enumerate}[label=(\roman*)]
\setlength{\itemsep}{0pt}
\item $[T]_{\beta\leftarrow\beta}$ is upper triangular;
\item $T\mathbf{v}_j\in\operatorname{span}(\mathbf{v}_1,\ldots,\mathbf{v}_j)$ for each $j=1,\ldots,n$;
\item $\operatorname{span}(\mathbf{v}_1,\ldots,\mathbf{v}_j)$ is invariant under $T$ for each $j=1,\ldots,n$.
\end{enumerate}
\begin{proof}
\begin{enumerate}[label=(\roman*).]
\setlength{\itemsep}{0pt}
\item[$(i)$]$\Rightarrow (ii)$
The $j$th column of $[T]_{\beta\leftarrow\beta}$ is the coordinate column of $T\mathbf{v}_j$ in the basis $\beta$. Since the matrix is upper triangular, this column has zero entries below position $j$, so $T\mathbf{v}_j\in\operatorname{span}(\mathbf{v}_1,\ldots,\mathbf{v}_j)$.
\item[$(ii)$]$\Rightarrow (i)$
If $T\mathbf{v}_j\in\operatorname{span}(\mathbf{v}_1,\ldots,\mathbf{v}_j)$ for every $j$, then the $j$th coordinate column of $T\mathbf{v}_j$ has zero entries below position $j$. Hence $[T]_{\beta\leftarrow\beta}$ is upper triangular.
\item[$(ii)$]$\Rightarrow (iii)$
If $\mathbf{u}\in\operatorname{span}(\mathbf{v}_1,\ldots,\mathbf{v}_j)$, then linearity gives $T\mathbf{u}$ as a linear combination of $T\mathbf{v}_1,\ldots,T\mathbf{v}_j$, all of which lie in $\operatorname{span}(\mathbf{v}_1,\ldots,\mathbf{v}_j)$. Hence this span is invariant.
\item[$(iii)$]$\Rightarrow (ii)$
Apply invariance of $\operatorname{span}(\mathbf{v}_1,\ldots,\mathbf{v}_j)$ to $\mathbf{v}_j$.
\end{enumerate}
\hlt{Structural Commentary}: Upper triangularity records an invariant flag.\\
\hlt{Key Insight}: Columns below the diagonal vanish exactly when $T\mathbf{v}_j$ uses only earlier basis vectors and itself.
\end{proof}
\end{theorem}

\begin{theorem} \label{thm:complex-upper-triangularization} \hlt{Complex Upper Triangularisation}\\
Let $\mathcal{V}$ be a finite-dimensional complex vector space and let $T\in\LL(\mathcal{V})$. Then $T$ has an upper-triangular matrix with respect to some basis of $\mathcal{V}$.
\begin{proof} Induct on $n=\dim(\mathcal{V})$. If $n\leq1$, the result is immediate.\\
Suppose $n>1$ and the result holds in smaller dimensions. Choose an eigenvector $\mathbf{v}_1$ of $T$ by Theorem \ref{thm:complex-operator-has-eigenvalue}, and set $\mathcal{U}=\operatorname{span}(\mathbf{v}_1)$. Then $\mathcal{U}$ is $T$-invariant, so the quotient operator $T/\mathcal{U}$ is defined on $\mathcal{V}/\mathcal{U}$, which has dimension $n-1$. By induction, choose a basis $(\mathbf{v}_2+\mathcal{U},\ldots,\mathbf{v}_n+\mathcal{U})$ of $\mathcal{V}/\mathcal{U}$ with respect to which $T/\mathcal{U}$ is upper triangular, and choose representatives $\mathbf{v}_2,\ldots,\mathbf{v}_n\in\mathcal{V}$. Then $(\mathbf{v}_1,\ldots,\mathbf{v}_n)$ is a basis of $\mathcal{V}$: any relation descends to the quotient basis, forcing all coefficients of $\mathbf{v}_2,\ldots,\mathbf{v}_n$ to be zero, and then the coefficient of $\mathbf{v}_1$ is zero.\\
Now $T\mathbf{v}_1\in\operatorname{span}(\mathbf{v}_1)$. For $j\geq2$, upper triangularity of $T/\mathcal{U}$ gives $(T/\mathcal{U})(\mathbf{v}_j+\mathcal{U}) \in\operatorname{span}(\mathbf{v}_2+\mathcal{U},\ldots,\mathbf{v}_j+\mathcal{U})$, so $T\mathbf{v}_j\in\operatorname{span}(\mathbf{v}_1,\ldots,\mathbf{v}_j)$. Hence Theorem \ref{thm:upper-triangular-criteria} gives an upper-triangular matrix for $T$ in the basis $(\mathbf{v}_1,\ldots,\mathbf{v}_n)$.\\
\hlt{Structural Commentary}: Lift an upper-triangular quotient basis above an eigenline.\\
\hlt{Key Insight}: The quotient builds the remaining invariant flag.
\end{proof}
\end{theorem}

\begin{proposition} \label{prop:upper-triangular-invertibility-criterion} \hlt{Invertibility from an Upper-Triangular Matrix}\\
Suppose $T\in\LL(\mathcal{V})$ has an upper-triangular matrix with diagonal entries $\lambda_1,\ldots,\lambda_n$ with respect to some basis of $\mathcal{V}$. Then $T$ is invertible if and only if $\lambda_j\neq0$ for every $j$.
\begin{proof}
Let $\beta=(\mathbf{v}_1,\ldots,\mathbf{v}_n)$ be such a basis.
\begin{enumerate}[label=\roman*.]
\setlength{\itemsep}{0pt}
\item[$(\Leftarrow)$]
If every $\lambda_j\neq0$, then upper-triangular matrix $[T]_{\beta\leftarrow\beta}$ is invertible by Theorem \ref{thm:inverse-upper-triangular}, so $T$ is invertible.
\item[$(\Rightarrow)$]
Suppose $T$ is invertible. If $\lambda_j=0$ for some $j$, set $\mathcal{U}_j=\operatorname{span}(\mathbf{v}_1,\ldots,\mathbf{v}_j)$. Upper triangularity gives $T(\mathcal{U}_j)\subseteq\mathcal{U}_{j-1}$ because the coefficient of $\mathbf{v}_j$ in $T\mathbf{v}_j$ is $\lambda_j=0$ and $T\mathbf{v}_i\in\mathcal{U}_{j-1}$ for $i<j$. Since $\dim(\mathcal{U}_j)>\dim(\mathcal{U}_{j-1})$, the restriction $T|_{\mathcal{U}_j}:\mathcal{U}_j\to\mathcal{U}_{j-1}$ is not injective. Hence $T$ has a nonzero kernel vector, contradicting invertibility. Thus every $\lambda_j\neq0$.
\end{enumerate}
\hlt{Structural Commentary}: Diagonal entries control each flag layer.\\
\hlt{Key Insight}: A zero diagonal entry forces a kernel vector.
\end{proof}
\end{proposition}

\begin{proposition} \label{prop:eigenvalues-of-upper-triangular-operator} \hlt{Eigenvalues from an Upper-Triangular Matrix}\\
Suppose $T\in\LL(\mathcal{V})$ has an upper-triangular matrix with diagonal entries $\lambda_1,\ldots,\lambda_n$ with respect to some basis of $\mathcal{V}$. Then the eigenvalues of $T$ are precisely $\lambda_1,\ldots,\lambda_n$.
\begin{proof}
For $\lambda\in\mathbb{F}$, the matrix of $T-\lambda I_{\mathcal{V}}$ in the same basis is upper triangular with diagonal entries $\lambda_1-\lambda,\ldots,\lambda_n-\lambda$. By Proposition \ref{prop:upper-triangular-invertibility-criterion}, $T-\lambda I_{\mathcal{V}}$ is not invertible if and only if $\lambda=\lambda_j$ for some $j$. By Theorem \ref{thm:equivalent-eigenvalue-conditions}, this is equivalent to $\lambda$ being an eigenvalue of $T$.\\
\hlt{Structural Commentary}: Scalar shifts preserve triangular form.\\
\hlt{Key Insight}: Eigenvalues occur exactly when one shifted diagonal entry becomes zero.
\end{proof}
\end{proposition}

\begin{perspective} \label{per:upper-triangular-viewpoint} \hlt{Why Triangular Form Matters}\\
Upper-triangular form is the first useful simplification beyond finding one eigenvector. It preserves enough of the operator to read eigenvalues from the diagonal while still allowing nontrivial coupling above the diagonal.
\end{perspective}
